\section{Connected localization of spatially canceled low kernels}
\label{lst-section}

This section proves a connected-activity norm estimate for the
canceled \emph{total} kernels. Its inputs are coherent complete-field
support families and exact infinite-lattice moment cancellations.
It also applies to explicitly periodized versions of those inputs,
with their lifted sums. It does not silently impose lifted moment
identities on an arbitrary finite torus kernel.

\subsection{Inputs and parameter order}
\label{lst-input-section}

Work in dimension four. A current coordinate consists of a site,
positive bond direction and color, with fixed internal dimension \(v\).
For \(n=2,3,4\), let \(K_n\) be a translation-invariant symmetric
\(n\)-linear kernel. Anchor its first site at \(x\), write the other
sites as \(x+r_2,\ldots,x+r_n\), and put
\(m=\sum_{j=2}^n|r_j|_1\). Internal indices are included in all sums.
Assume
\begin{equation}
 \sup_{\text{first internal index}}
 \sum_{\text{other indices},r_2,\ldots,r_n}
       e^{\sigma m}|K_n(0,r_2,\ldots,r_n)|
       \le \epsilon b^{-n}.
 \label{lst-kernel-bound}
\end{equation}
The constants \(\sigma,\epsilon\) can incorporate the fixed old
grouping \(J\). Fixed-order Cauchy and a connected activity bound
supply this estimate with a smaller spatial exponent.

Use tail-site coordinates. Assume componentwise
\begin{align}
 \sum_r r^\gamma K_2(r)&=0
                          &&(|\gamma|\le3),\notag\\
 \sum_{r,s}K_3(r,s)&=0,\qquad
 \sum_{r,s}r_\rho K_3(r,s)=
 \sum_{r,s}s_\rho K_3(r,s)=0
                          &&(\rho=1,\ldots,4),\notag\\
 \sum_{r,s,t}K_4(r,s,t)&=0.
 \label{lst-moments}
\end{align}
These are consequences of the stronger vanishing orders in
Theorem~\ref{tkw-theorem}; analytic conversion from midpoint
to tail coordinates preserves those orders.

Fix an integer endpoint ratio \(L\), put \(h=L^{-1}\), and let
\(a(w)\) be an analytic current-angle background on the radius-\(b\)
final-field ball. Current sites have physical spacing \(h\).
Assume coherent complete-field strict families: on every fixed
enlarged final-unit collar rooted at \(i\),
\begin{equation}
 a|_i=\sum_{X\ne\varnothing}a_{i,X}(w_X),\qquad
 a_{i,X}(0)=0,\qquad
 \sup_i\sum_X e^{\kappa\tau(\{i\}\cup X)+A|X|}
           \|a_{i,X}\|_{C^{1,\alpha};b}
       \le C_* b h .
 \label{lst-background}
\end{equation}
The derivatives in this norm are physical. Coefficients on
overlapping collars restrict one global assignment. All maxima,
first derivatives and first-gradient H\"older quotients are
inside the support sum. A fixed collar contains any coordinate
paths whose current length is at most \(L\); enlarging it changes
only fixed constants and possibly decreases \(\kappa\).
These are the complete-family inputs of the supported pre-gauge
background construction, not bounds merely on the summed field.

Let \(P_n(a)=K_n(a,\ldots,a)/n!\), fix \(a_0>0\), and set
\(c=a_0+1\). Choose \(J\) once with
\begin{equation}
 \frac{2c}{J}\le\kappa,\qquad A\ge4c .
 \label{lst-budgets}
\end{equation}

\begin{theorem}[Low-kernel spatial transport]
\label{lst-theorem}
There are \(C,L_0\), depending on
\(\sigma,J,C_*,a_0,\alpha\) and the fixed internal dimension,
but not on \(L,b\), volume or preceding refinement, such that
for \(L\ge L_0\) there are connected activities on final \(J\)-cells
with
\begin{equation}
 \sum_H B_H(w_H)=\sum_{n=2}^4P_n(a(w)),\qquad
 \sup_I\sum_{H\ni I}e^{a_0|H|}\|B_H\|_b
                      \le C\epsilon L^{-\alpha}.
 \label{lst-conclusion}
\end{equation}
The construction is linear in the total kernels.
Individual \(B_H\) need not be gauge invariant.
The scalar identity is vacuum-normalized and absolutely
convergent for compactly supported final fields.
\end{theorem}

The constant is fixed after \(J\) and before \(L\).
One can therefore set \(L=R^q\) with fixed \(R,J\) and
choose a finite \(q\) afterward. This statement does not
identify \(q\) successive quantum integrations with one
direct fluctuation integral.

\subsection{A valid quadratic difference factorization}
\label{lst-quadratic-section}

The matrix kernel need not be symmetric at each displacement,
even though Hessian reciprocity gives
\(K_2(r)^T=K_2(-r)\). We use a Laurent factorization that
does not require pointwise symmetry. Set
\[
 F(z)=\sum_r K_2(r)z^r,\qquad
 d_\rho^+(z)=z_\rho-1,\qquad
 d_\rho^-(z)=z_\rho^{-1}-1 .
\]
The exponential bound makes \(F\) holomorphic on a product
annulus. Equation~\eqref{lst-moments} says its jets through
degree three at \(z=(1,\ldots,1)\) vanish.

For a Laurent function define the ordered divided difference
\begin{equation}
 {\cal B}_jF(z)=
 \frac{F(1,\ldots,1,z_j,\ldots,z_4)
        -F(1,\ldots,1,1,z_{j+1},\ldots,z_4)}
      {z_j-1}.
 \label{lst-divided-difference}
\end{equation}
Its singularity at \(z_j=1\) is removable, and
\[
 F(z)-F(1)=\sum_j(z_j-1){\cal B}_jF(z).
\]
For a monomial \(z^r\), the quotient in coordinate \(j\)
has \(|r_j|\) monomials, each at absolute exponent length
at most \(|r|_1\), including negative exponents. Therefore,
for \(0<\sigma'<\sigma\), the combined absolute row Wiener
norm obeys
\begin{equation}
 \sup_u\sum_{j,v,r}e^{\sigma'|r|_1}
       |({\cal B}_jF)_{uv}(r)|
 \le\frac1{e(\sigma-\sigma')}\,
             \sup_u\sum_{v,r}e^{\sigma|r|_1}|F_{uv}(r)|.
 \label{lst-Wiener-quotient}
\end{equation}
Here the \(j\) sum is inside the row supremum. The bound
uses \(\sum_j|r_j|=|r|_1\) and
\(\sup_{t\ge0}t e^{-(\sigma-\sigma')t}
 =[e(\sigma-\sigma')]^{-1}\).
The local Taylor series also shows that
\({\cal B}_j\) lowers a vanishing order by at most one.

Apply the identity twice, using exponents
\(\sigma,3\sigma/4,\sigma/2\). Put
\[
 H_{\rho\tau}={\cal B}_\tau{\cal B}_\rho F,\qquad
 \widetilde H_{\rho\tau}(z)=-z_\rho H_{\rho\tau}(z).
\]
Since \(d_\rho^-(-z_\rho)=d_\rho^+\), this gives
\begin{equation}
 F(z)=\sum_{\rho,\tau}
       d_\rho^-(z)\widetilde H_{\rho\tau}(z)d_\tau^+(z),
 \qquad \widetilde H_{\rho\tau}(1)=0.
 \label{lst-factorization}
\end{equation}
In fact the jets of \(\widetilde H_{\rho\tau}\) through
degree one vanish. Its combined kernel bound is
\begin{equation}
 \sup_u\sum_{\rho,\tau,v,r}
     e^{(\sigma/2)|r|_1}
           |(\widetilde H_{\rho\tau})_{uv}(r)|
 \le \frac{16e^{\sigma/2}}{e^2\sigma^2}
                           \epsilon b^{-2}.
 \label{lst-H-bound}
\end{equation}
The factor \(e^{\sigma/2}\) pays the unit shift. No symmetry
of the individual \(\widetilde H_{\rho\tau}\) is asserted.

Let \(D_\rho a(x)=a(x+e_\rho)-a(x)\), the unscaled current
difference. Integration by parts transfers one difference
to each field. The zero sum in \eqref{lst-factorization}
then gives the exact identity
\begin{equation}
 P_2(a)=\frac12\sum_{x,r,\rho,\tau}
 [D_\rho a(x)]^T\widetilde H_{\rho\tau}(r)
            [D_\tau a(x+r)-D_\tau a(x)].
 \label{lst-quadratic-identity}
\end{equation}
The sign follows because \(d_\rho^-\) is the bilinear
adjoint of \(d_\rho^+\). Compactly supported fields give
the identity directly; the absolute majorants below
extend it to the substituted fields. No positivity or
second field derivative is used.

\subsection{Cubic and quartic cancellation}
\label{lst-higher-algebra}

Put
\begin{equation}
 \ell_r(x)=\sum_\rho r_\rho D_\rho a(x),\qquad
 R_r(x)=a(x+r)-a(x)-\ell_r(x).
 \label{lst-Taylor}
\end{equation}
In the anchored cubic expression
\[
 P_3(a)=\frac16\sum_{x,r,s}
                K_3(r,s)[a(x),a(x+r),a(x+s)],
\]
replace each shifted factor by \(a(x)+\ell+R\).
The zeroth moment cancels the all-anchor term. The
first moments cancel \(a,\ell,a\) and \(a,a,\ell\).
Six terms remain: the two with one \(R\) and no \(\ell\),
the \(\ell,\ell\) term, the two mixed \(\ell,R\) terms,
and the \(R,R\) term, always with unshifted first factor.

For the quartic expression subtract the all-anchor product
using its zeroth moment, then telescope the three shifted
factors. Each of three terms has one difference
\(a(x+r_j)-a(x)\) and three endpoint or anchor factors.
All cancellations are made in the total kernels before
triangle inequalities or final carrier assignment.

\subsection{Near-displacement coefficient bounds}
\label{lst-near-section}

Split these exact residual sums into \(m\le L\) and \(m>L\).
Neither restricted kernel is claimed to retain zero moments.
For the quadratic identity use \(m=|r|_1\).
In the near case all evaluations and coordinate paths lie
in a fixed final-unit collar. Coherence in
\eqref{lst-background} gives rooted coefficient-family bounds
\begin{align}
 |a|&\le Cbh,\qquad
 |\ell_r|\le Cbh^2|r|_1,\notag\\
 |a(x+r)-a(x)|&\le Cbh^2|r|_1,\qquad
 |R_r|\le Cbh^{2+\alpha}|r|_1^{1+\alpha}.
 \label{lst-near-fields}
\end{align}
For a negative coordinate traversal, the positive difference
is evaluated at the lower endpoint with the traversal sign.
Subtracting its value at the anchor and using the physical
gradient H\"older bound proves the last estimate.
Each inequality is for the weighted sum of coefficient
norms, with its suprema and H\"older quotients inside that
sum. It is not deduced from a bound on the scalar sum.

All near factors can be rooted at the anchor cell.
Multiplication glues rooted trees and support unions, so
their weights are submultiplicative. In the quadratic
identity,
\[
 |D_\rho a(x)|\le Cbh^2,\qquad
 |D_\tau a(x+r)-D_\tau a(x)|
                      \le Cbh^{2+\alpha}|r|_1^\alpha.
\]
Its product is therefore
\begin{equation}
 Cb^2h^{4+\alpha}|r|_1^\alpha.
 \label{lst-near-quadratic}
\end{equation}
Sum it with \eqref{lst-H-bound}. The six cubic products
are bounded by
\begin{equation}
 Cb^3h^{4+\alpha}(1+m)^6.
 \label{lst-near-cubic}
\end{equation}
The leading \(a,a,R\) term has that power of \(h\);
\(a,\ell,\ell\) has \(h^5\), and the others have still
higher powers. The quartic telescoping is bounded by
\begin{equation}
 Cb^4h^5(1+m).
 \label{lst-near-quartic}
\end{equation}
The weighted kernel sums cancel \(b^n\) and pay all
fixed displacement moments. Per original current anchor,
the three bounds are respectively
\(C\epsilon h^{4+\alpha}\),
\(C\epsilon h^{4+\alpha}\), and \(C\epsilon h^5\).

\subsection{Connected carriers and the anchor count}
\label{lst-carriers-section}

For each product coefficient keep its final-unit anchor
\(i_0\) and the endpoint roots actually used. In a near
term the complete collar permits the common root \(i_0\).
In a far term use the endpoint roots and any neighboring
roots needed for \(Da\). Their number is at most a fixed
\(n_0\), independent of \(L\) and the displacement.

Project the finite root set and, for each field-factor
occurrence, the terminal set consisting of its root
and variable support \(X\) to final \(J\)-cells.
For each projected terminal set choose a deterministic
short macro connecting graph supplied by the projection
lemma. Take the connected union of these macro graphs
as \(H\), including the anchor. The fine Steiner
lengths below bound those chosen macro graphs; the
carrier need not contain the full projection of an
arbitrary fine connecting tree. The projection bound gives
\begin{align}
 |H|&\le1+T,\notag\\
 T&\le\frac2J\left[
       \tau(\text{roots})+
       \sum_{\mathrm{occ}}\tau(\{\text{root}\}\cup X)\right]
         +4\left[n_0+\sum_{\mathrm{occ}}|X|\right].
 \label{lst-tree-bound}
\end{align}
Let \(j(x)\) be the final \(J\)-cell containing \(x/L\).
For \(z\in H\), its graph distance to \(j(x)\) is at
most \(T\). Hence
\begin{equation}
 e^{a_0|H|}1_{\{z\in H\}}
 \le e^{a_0}e^{-\operatorname{dist}(z,j(x))}e^{cT}.
 \label{lst-reroot}
\end{equation}
The budgets \eqref{lst-budgets} pay every attached
tree and support-cardinality factor. The fixed
\(e^{4cn_0}\) is absorbed into \(C\); near root-joining
trees have bounded length.

There are exactly \((JL)^4\) current anchor sites
per final \(J\)-cell, with the same centered-boundary
convention throughout. Therefore
\begin{equation}
 \sum_{x\in\mathbb Z^4}
 e^{-\operatorname{dist}(z,j(x))}
 \le (JL)^4 S_1,\qquad
 S_1=\sum_{y\in\mathbb Z^4}e^{-|y|_1}.
 \label{lst-anchor-count}
\end{equation}
The first internal-index count is the fixed \(v\).
There is no additional count of all output coordinates.
The near contributions are consequently
\(CJ^4\epsilon L^{-\alpha}\) for degrees two and
three, and \(CJ^4\epsilon L^{-1}\) for degree four.

\subsection{Far displacements and convergence}
\label{lst-far-section}

For \(m>L\), use the same exact residual formulas but
root each endpoint separately. The cubic and quartic
factors obey
\[
 |a|\le Cbh,\qquad
 |\ell_r|\le Cbh^2m,\qquad
 |R_r|\le Cbh(1+hm).
\]
Their products are bounded by
\(Cb^3h^3(1+m)^2\) and \(Cb^4h^4\).
For the quadratic expression both first differences
are bounded by \(Cbh^2\) at their own roots, so its
product is bounded by \(Cb^2h^4\).

The finite endpoint-root tree has
\[
 \tau(\text{roots})\le C_d(m/L+1).
\]
After \eqref{lst-reroot} this costs at most
\(C\exp(C'm/(JL))\). Use the common kernel exponent
\(\sigma_*=\sigma/2\), supplied by \eqref{lst-H-bound}
and also available for \(K_3,K_4\). Choose \(L_0\)
so that \(C'/(JL)\le\sigma_*/4\). One half of
\(\sigma_*\) pays the extra tree weight and the fixed
polynomial, leaving an exponentially small tail.
After \eqref{lst-anchor-count} the total far bound is
\begin{equation}
 CJ^4\epsilon(L+2)e^{-\sigma_*L/4}
                 \le C'\epsilon L^{-\alpha},
       \qquad L\ge L_0 .
 \label{lst-far-bound}
\end{equation}
All constants here are fixed before \(L\).

The absolute majorants justify regrouping by connected
\(H\) and give holomorphic coefficients on the same
open \(b\)-ball, with their bounded supremum norm.
For compactly supported \(w\), each nonzero field-factor
coefficient has variable support meeting that compact
set, since it vanishes at zero. Its final carrier
contains that support. The anchored bound therefore
implies convergence of the total sum, bounded by a
constant times the number of final \(J\)-cells meeting
the support. For that compact support \(K\), every
nonzero field coefficient also has
\(\tau(\{i\}\cup X)\ge\operatorname{dist}(i,K)\).
Equation~\eqref{lst-background} therefore gives
exponential spatial decay of \(a(w)\) and its first
differences. The original quadratic kernel sum is
absolutely convergent as well, so truncation and
exhaustion justify its integration by parts without
a boundary term. The near and far
bounds prove Theorem~\ref{lst-theorem}.

\subsection{Scope of the relocalization}
\label{lst-scope}

The result is a bounded connected localization of
the canceled low-degree \emph{total} functional.
It is stronger than extensive scalar power counting
and needs no invariant individual carriers. For a
residual analytic activity \(E-\beta S\), its
coefficientwise degree-\(\ge5\) part can be transported
by the previous high-jet estimate, while the total
degrees two through four use the present construction.
The two localizations may change the representation.

This section does not prove canonical finite-torus
wrapping control, an infinite-volume quantum measure,
equality of sequential and direct cutoff integrals,
boundedness of a complete quantum map on a quotient
of representations, or a coupling trajectory.
The classical scalar semigroup permits the choice
\(L=R^q\) for a gauge-invariant total, but removes
none of those separate interfaces.
